\documentclass{ieeeaccess}
\usepackage{cite}
\usepackage{amsmath,amssymb,amsfonts}
\usepackage{graphicx}
\usepackage{textcomp}
\usepackage{alltt}
\usepackage{array}
\usepackage{tabularx}
\renewcommand{\thetable}{S\arabic{table}}
\renewcommand{\thefigure}{S\arabic{figure}}
\begin{document}
\history{Supplementary material.}
\doi{10.1109/ACCESS.2026.0429000}
\title{Supplement S1 --- Protocol and provenance record}
\author{\uppercase{Jay Suresh Nirmal}\authorrefmark{1}}
\address[1]{Independent Researcher, Boston, MA, USA}
\markboth{Nirmal: Supplement S1}{Nirmal: Supplement S1}

\begin{abstract}
Supplementary material to "Measuring the Safe Drain Boundary for Durable Backlog Recovery Under a Live Latency Objective". Every item here is cited from the article, which states each item's consequence; this document gives the full record. Text moved from the article is reproduced unchanged.
\end{abstract}

\titlepgskip=-15pt
\maketitle

\section*{S1-A. Provenance (full text of the article's former \S{}IV-G)}

Early in the project 137 of 140 run records carried the string \texttt{''unknown''} in place of a commit hash, because the function resolving it swallowed its error on an unborn HEAD. The runner now refuses to start without a resolvable commit. That guard covers the runner's own records --- every boundary run carries a commit, a branch, a dirty flag and a start time. It does not cover the standalone saturation tool that measures plateaus, which stamps no provenance of any kind; \S{}V states what that costs and \S{}IX carries it. No run reported in this paper lacks a commit.

One provenance gap remains and is reported rather than closed. The search driver records its results but not its own invocation, so the arguments that produced a cell's probe set are recoverable only by replaying the committed planner against the observed classifications. All seven cells reconstruct exactly, each with no supplied ceiling. One required care: in E1 c10/C1 two SAFE points lie below the floor the descent located, and they carry a run-id prefix identifying them as probes from the replication campaign rather than the original search, which the artefact does record. Replaying the planner on the original campaign's rates alone reproduces that cell's probe set exactly.

What is not recoverable is the command line itself. The gap is therefore in the reproducibility of probe \textit{selection}, not in the measurements: for every cell, all probed rates, all repetitions, all classifications and the terminal 5 rps bracket are retained and used, no run carries an exclusion reason, and none is absent from its cell's analysis --- the search driver raises rather than continuing if a run is invalid, so a silently dropped run is not a state it can reach. Adding an argument echo would prevent a recurrence and would recover nothing retrospectively; the driver is left unaltered, because changing a frozen instrument after its campaigns have closed is the practice this section exists to avoid.

\section*{S1-B. Amendments (full text of the article's former \S{}IV-H, with Table S1)}

All ten amendments and record corrections are listed in Table~\ref{tab:s1-amendments} with their commits; the article glosses the six that later sections cite.

\begin{table*}[!t]
\caption[Amendments to the pre-registration, the A7 registered re-analysis, the A8 registered replication and the two record corrections.]{Amendments to the pre-registration, the A7 registered re-analysis, the A8 registered replication and the two record corrections. The two right-hand columns separate registration \textit{before the data existed} from registration \textit{before the analysis was run}; they are different claims and only the first is prospective registration. They do not apply to a record correction, which governs neither, and those rows read N/A rather than being forced into a category they do not belong to.}
\label{tab:s1-amendments}
\centering
\footnotesize
\begin{tabularx}{\textwidth}{lll>{\raggedright\arraybackslash}X>{\raggedright\arraybackslash}X>{\raggedright\arraybackslash}X}
\hline
{\bfseries \#} & {\bfseries Date} & {\bfseries Commit} & {\bfseries Change} & {\bfseries Predates governed data} & {\bfseries Predates governed analysis} \\ \hline
A1 & 2026-09-11 & \texttt{e5bf303} & Interval's upper end is the first NON-SAFE point, which may be MARGINAL, not the first UNSAFE point & Yes & Yes \\
A2 & 2026-09-11 & \texttt{1fe9de3} & A non-SAFE anchor extends the search downward instead of aborting for a hand-chosen anchor & Yes & Yes \\
A3 & 2026-09-11 & \texttt{699e105}, \texttt{25583c7} & Live injector moves from a tick-dropping timer to lanes & Yes & Yes \\
A4 & 2026-09-11 & \texttt{d5ea89d}, \texttt{ddd428d} & Achieved recovery rate measured over the delivery span, not the drain window & {\bfseries Partly} --- the first boundary's traces (c10/C0) already existed and were recomputed; registered before the remaining E1 cells & Yes \\
A5 & 2026-09-12 & \texttt{fe36734} & Leading-indicator noise scale changed from mean to median & {\bfseries No} & {\bfseries No} --- chosen with the numbers in view \\
A6 & 2026-09-13 & \texttt{c3aee75} & A4 is not valid at collapsed points; utilisation is not reported there & {\bfseries No} & {\bfseries No} \\
A7 & 2026-09-14 & \texttt{939902d} & \textit{Registered re-analysis, not an amendment.} The frozen leading-indicator criterion and statistic, re-run over the corrected-harness corpus & {\bfseries No} --- that corpus already existed & {\bfseries Yes} --- criterion, statistic and prediction all frozen before the re-analysis \\
A8 & 2026-09-14 & \texttt{590d1cc} & \textit{Registered replication, not an amendment.} Repeatability of the two corrected saturation plateaus: ten 60-second windows per arm, median convention, with both arms' readings fixed in advance. No registered constant, no plateau estimator and no analysis path is changed by it & {\bfseries Yes} --- the twenty windows were collected afterwards & {\bfseries Yes} \\
A9 & 2026-09-19 & \texttt{9caf476} & \textit{Record correction, not an amendment.} A8's short-arm reading states that neither plateau figure was itself registered as a plateau; \texttt{67c448b} tabulates both and gates the boundary read on them, so that statement is wrong. No A8 criterion, datum or result depends on it & {\bfseries N/A} --- governs no data & {\bfseries N/A} --- governs no analysis \\
A10 & 2026-09-19 & \texttt{51026cc} & \textit{Record correction, not an amendment.} A6 reports the E2d collapse factor as 10.3$\times$, computed from a midpoint numerator against a SAFE-only denominator; matched on the SAFE side it is 10.5$\times$. The denominator, the direction and A6's verdict are unchanged & {\bfseries N/A} --- governs no data & {\bfseries N/A} --- governs no analysis \\
\hline
\end{tabularx}
\end{table*}

Amendments to the pre-registration, the A7 registered re-analysis, the A8 registered replication and the two record corrections. The two right-hand columns separate registration \textit{before the data existed} from registration \textit{before the analysis was run}; they are different claims and only the first is prospective registration. They do not apply to a record correction, which governs neither, and those rows read N/A rather than being forced into a category they do not belong to.

A7 is a registered re-analysis of an already-collected corpus, not a prospective registration, and the two columns are separated so that the distinction cannot be read the other way.

A8 is neither an amendment nor a re-analysis. It changes no rule and revisits no existing data: it registers a \textit{new} measurement --- the window-to-window repeatability of two plateaus the campaign had measured once each --- before that measurement was taken, and fixes in advance what each possible outcome would mean, including the outcome that would have removed a prospective claim \S{}V currently makes. It is prospective in the same strict sense as A1 through A3, and appears last only because it postdates the measurement campaign it bears on. Its result is reported in \S{}V-D.

A9 and A10 are a third kind again: record corrections. Each fixes a factual misstatement in the text of an earlier registration, and neither changes a rule, an estimator, a datum or a conclusion. They exist because this registration is immutable once its commit is cited, so a correction takes the form of a dated addendum with the original left standing --- and because applying that rule only when a correction is convenient would cost more than the addenda do. A9's error understated what had been registered; A10's makes a reported factor 2\% larger. The second direction warrants more scrutiny than the first, and \S{}IX gives both in full.

A1 through A3 were registered before any data governed by this registration were collected --- an earlier phase of the project had produced data, which is what motivated the protocol --- and require no further comment. A8 is described above; the remaining four do.

A4 replaced an estimator, not an instrument. Chronologically it was registered during E1, after the traces for the first boundary (c10/C0) already existed and before the remaining E1 cells were collected; it was then applied retrospectively to those first traces. The original measure divided an exact backlog by a window padded with the drain detector's confirmation tail, which injected run-to-run noise the rate limiter had not produced. Measuring instead over the span the recovery traffic actually occupied removes that tail. It was applied to already-recorded traces and required no run to be repeated.

A5 is the one amendment in this document made with the numbers already in view, and it changes an answer. The registration says so in those terms and documents it at length precisely because it is the most attackable step in the leading-indicator analysis. It is reported here on the same basis: a reader should be able to attack it with the evidence in front of them.

A7 attempts to repair that, and does not succeed. A5 froze the statistic on 12 September at 07:31 UTC, having been chosen on the E1 corpus collected five hours earlier --- and the leading-indicator analysis reads only that corpus. A7 therefore registered, before running anything, a re-analysis of the same frozen criterion and the same frozen statistic over the corrected-harness corpus, which was collected 34 hours after A5 froze and played no part in choosing it, together with the prediction to be tested. The registration preceded the analysis by commit timestamp.

The prediction did not replicate. The timeout component held cleanly. The ordering component did not: queue depth and live p99 first move at the same probed point in both corrected cells. The registered limitation is what decides the interpretation. E1's safe grids are sampled at 5--10 rps where they are finest; the corrected cells at 60 and 100 rps. A lead of the 10--25 rps that E1 reported is below the corrected corpus's resolution by construction, so that corpus could not have replicated the ordering in either direction. The honest statement is the stronger one: the ordering is unreplicated, and the only independent corpus available was incapable of adjudicating it. \S{}VII reports it as exploratory on those grounds and on A5's chronology independently.

A6 records that A4 is valid at safe points but over-reads at collapsed ones, where the recovery traffic's occupied span is not the drain. Eighteen of twenty unsafe points reported achieved rates above their own saturation plateau --- a count, not a magnitude, and therefore not a judgement about resolution. The defect was caught by a cross-check of the estimator against the plateau rather than by inspection, and utilisation is no longer reported at collapsed points. Its consequences for the reported intervals are set out in \S{}IX.

\section*{S1-C. Reproducibility (full text of the article's former \S{}IV-I)}

The harness, the pre-registration with all six amendments, the A7 registered re-analysis, the A8 registered replication and the A9 and A10 record corrections, every run record and per-request trace --- the traces are archived with the dataset rather than held in the repository, so delivery-span quantities are checkable against their artefacts from the repository alone but re-derivable from raw observations only with the deposit --- the analysis code, the figure generators and the regression fixture are archived under DOI \texttt{10.5281/zenodo.22761131}. Every reported artefact regenerates byte-identically from the committed data by running its own script, and a manifest records a SHA-256 for each file together with the commit the package was built from.

That property was verified rather than assumed, and verifying it found a counter-example. The data artefacts reproduced byte-identically throughout; the figures did not, because the plotting library stamps a wall-clock creation time into each file, so two runs over unchanged inputs produced different bytes with identical drawn content. The metadata is now suppressed and determinism is verified over consecutive runs, with the rendered marks confirmed unchanged against the committed versions. The claim as stated held for the artefacts it had been checked against and not for the ones it had not, which is the ordinary way such claims fail.

\section*{S1-D. The capacity model, pictured (the article's former Fig. 2)}

Fig.~\ref{fig:capacity-model} pictures the relation equations \eqref{eq:4}--\eqref{eq:6} state algebraically: the harness staffs its pool on the assumption that a worker turns a request round in exactly the emulated service time, each worker in fact pays a fixed additional cost per request, and because that cost is additive rather than proportional it consumes a larger share of a short service time than a long one. The figure carries no measured result, which is why the article states the relation in equations and keeps the picture here.

\begin{figure*}[!t]
\centering
\includegraphics[width=\textwidth]{F2-capacity-model.pdf}
\caption[Where the unaccounted cost enters.]{Where the unaccounted cost enters. The harness sizes the dependency as c = $\lceil$C\_config$\cdot$S$\rceil$ workers on the assumption that each turns a request round in exactly S, which makes capacity equal to C\_config by construction. A worker in fact pays a fixed additional 0.463 ms per request, almost all of it the operating system sleeping longer than asked. C\_model is therefore c/(S + $\delta$), as in (5), and because the cost is additive rather than proportional it consumes a larger share of a short service time than a long one.}
\label{fig:capacity-model}
\end{figure*}

\section*{S1-E. Per-cell resolution in full (the article's Table 2, all columns)}

The article's resolution table is compacted to the three quantities its argument uses: the boundary bracket, the normalised per-cell resolution and the safe utilisation. Table~\ref{tab:resolution-full} is the same seven cells with every column the generator produces, so that a reader checking the resolution discipline against the artefacts has the configured and measured capacities, the raw bisection step, the bracket width, the repetition count, the replicate spread and the measured capacity's own range in one place.

\begin{table*}[!t]
\caption[Per-cell resolution of the boundary estimate, in full: the article's compact table gives the bracket, the normalised resolution and the safe utilisation; this one adds the configured and measured capacities, the raw bisection step, the bracket width, the repetition count, the replicate spread and the measured capacity's own range.]{Per-cell resolution of the boundary estimate, in full: the article's compact table gives the bracket, the normalised resolution and the safe utilisation; this one adds the configured and measured capacities, the raw bisection step, the bracket width, the repetition count, the replicate spread and the measured capacity's own range.}
\label{tab:resolution-full}
\centering
\footnotesize
\begin{tabularx}{\textwidth}{lrrrrlrr>{\raggedright\arraybackslash}Xr>{\raggedright\arraybackslash}X}
\hline
{\bfseries cell} & {\bfseries C\_config} & {\bfseries C\_measured} & {\bfseries step (rps)} & {\bfseries step / C\_measured} & {\bfseries rate bracket (rps)} & {\bfseries width (rps)} & {\bfseries n} & {\bfseries replicate spread (rps)} & {\bfseries $\rho$\_eff,safe} & {\bfseries C\_measured range (rps)} \\ \hline
E1 c10/C0 & 2000 & 1828.6 & 5 & 0.0027 & [825, 830] & 5 & 15 & 0.2 & 0.998 & 1828.0 -- 1829.9 (n = 6) \\
E1 c10/C1 & 1400 & 1280.9 & 5 & 0.0039 & [275, 280] & 5 & 3 & 0.0 & 0.995 & 1279.1 -- 1282.0 (n = 6) \\
E1 c50/C0 & 2000 & 1964.3 & 5 & 0.0025 & [975, 980] & 5 & 15 & 1.2 & 1.000 & 1962.7 -- 1965.3 (n = 6) \\
E1 c50/C1 & 1400 & 1375.7 & 5 & 0.0036 & [370, 375] & 5 & 3 & 0.0 & 0.996 & 1375.0 -- 1378.3 (n = 6) \\
E2 c10@Q2500 & 2000 & 1827.1 & 5 & 0.0027 & [825, 830] & 5 & 12 & 0.2 & 0.999 & 1825.5 -- 1829.5 (n = 15) \\
E2 c50@Q500 & 2000 & 1964.8 & 5 & 0.0025 & [975, 980] & 5 & 12 & 1.2 & 1.000 & 1963.7 -- 1965.7 (n = 15) \\
E2b C=400 & 400 & 392.7 & 5 & 0.0127 & [190, 195] & 5 & 12 & 0.0 & 0.99 & 392.5 -- 393.0 (n = 6) \\
\hline
\end{tabularx}
\end{table*}


\EOD
\end{document}
